% =====================================================================
% Problem statement -- carries over the content of the approved topic
% proposal. Please do not change the structure.
% =====================================================================
\vorspannkapitel{Problemstellung}
\abschnittsautor{\alleautoren}

\vorspannabschnitt{Ausgangslage}

Beschreibt hier in wenigen Sätzen, welches Problem im Ausgangszustand besteht,
wen es betrifft und warum eine technische Lösung sinnvoll ist.

\vorspannabschnitt{Untersuchungsanliegen der individuellen Themenstellungen}

\vorspannunterpunkt{1) Teilbereich A}
Was untersucht dieses Teilprojekt? Formuliert die Fragestellung, nicht die
Lösung.

\vorspannunterpunkt{2) Teilbereich B}
\dots

\vorspannunterpunkt{3) Teilbereich C}
\dots

\vorspannabschnitt{Zielsetzung}

Ein bis zwei Sätze: Was soll am Ende der Arbeit vorliegen?

\vorspannabschnitt{Geplantes Ergebnis der individuellen Themenstellungen}

\vorspannunterpunkt{1) Teilbereich A}
Welches konkrete, überprüfbare Ergebnis liefert dieses Teilprojekt ab?

\vorspannunterpunkt{2) Teilbereich B}
\dots

\vorspannunterpunkt{3) Teilbereich C}
\dots

\vorspannabschnitt{Meilensteine}

% Deliberately not a table float: the milestones belong directly under the
% heading, not somewhere else in the document.
\begin{tabular}{@{}p{0.68\linewidth}p{0.24\linewidth}@{}}
    \toprule
    \bfseries Meilenstein & \bfseries Datum \\
    \midrule
    Analyse der Ausgangslage und Anforderungen & 30.09.2026 \\
    Erstellung der Konzepte                    & 31.10.2026 \\
    Implementierung der Systemkomponenten      & 30.11.2026 \\
    Test der Systemkomponenten                 & 15.01.2027 \\
    Integration des Gesamtsystems              & 15.02.2027 \\
    Verfassen der Dokumentation                & 15.03.2027 \\
  Abgabe der Diplomarbeit                    & 01.04.2027 \\
  \bottomrule
\end{tabular}
